\scanpage{26}
en associant à $(g,(\alpha_i)_{i\in\{0,1,\infty\}})$ le système
\[
\deleted{\text{de}}\quad p=1,\quad(g_i)=\deleted{\unsure{(l_i^{\alpha_i})}}\deleted{\unsure{(l_i^{\alpha_i}g^{-1})}}(gl_i^{\alpha_i})
\]
Il faudrait expliciter les opérations des éléments $\sigma_i$ ou $\varepsilon_i$ \deleted{sur} $(\in S\widehat{\widehat{\mathcal E}}^!)$ sur $\deleted{\unsure{u}}\tilde u\in S\widehat{\widehat{\mathcal E}}^!$, par \deleted{\unsure{leurs}} conjugaison, et en particulier de $\rho$. On trouve que \deleted{\unsure{$p$ (bien sûr)}} les $\gamma_i$ sont ``covariants'' par $\rho$, i.e.
\[
\gamma_i(\rho\tilde u\rho^{-1})=\gamma_{\unsure{i+1|i-1}}(\unsure{\bar u|\tilde u})
\]
\deleted{si $\tilde u$ est défini par}
\margin{\uline{NB} Il y a lieu d'introduire un groupe $S\unsure{\hat{\mathcal E}|\widehat{\widehat{\mathcal E}}}$ extension de $\mathcal E$ par $\hat{\mathbb Z}^{\{0,1,\infty\}}$ (déduite d'une extension canonique de $\Gamma$ par $\mathbb Z^{\{0,1,\infty\}}$) et d'introduire \deleted{\unsure{des}} des relèvements des $\sigma_i,\rho$ par des éléments de $S\unsure{\hat{\mathcal E}|\widehat{\widehat{\mathcal E}}}$\ldots}
\mspage{25}
